% !TeX program = xelatex
% Báo cáo tự chứa; biên dịch bằng XeLaTeX (khuyến nghị) hoặc pdfLaTeX có VnTeX.
% Chỉ mô tả chức năng có trong source; hình minh họa không bổ sung thuật toán mới.
\documentclass[10pt,a4paper]{article}
\usepackage{iftex}
\ifPDFTeX
  \usepackage[utf8]{inputenc}
  \usepackage[T5]{fontenc}
  \usepackage[vietnamese]{babel}
\else
  \usepackage{fontspec}
  \IfFontExistsTF{Latin Modern Roman}
    {\setmainfont{Latin Modern Roman}}{\setmainfont{Times New Roman}}
  \IfFontExistsTF{Latin Modern Mono}
    {\setmonofont{Latin Modern Mono}[Scale=MatchLowercase]}
    {\setmonofont{Consolas}[Scale=MatchLowercase]}
\fi
\usepackage[margin=1.7cm]{geometry}
\usepackage{amsmath,booktabs,array,tabularx,enumitem,fancyvrb}
\usepackage{tikz}
\usetikzlibrary{arrows.meta,positioning,calc}
\usepackage[unicode,hidelinks]{hyperref}
\setlength{\parindent}{0pt}
\setlength{\parskip}{4pt}
\setlist[itemize]{leftmargin=1.3em,itemsep=2pt,topsep=3pt}
\renewcommand{\arraystretch}{1.13}
\newcolumntype{Y}{>{\raggedright\arraybackslash}X}
\newcommand{\shared}{\texttt{shared}}
\newcommand{\percam}{\texttt{per\_camera}}
\DefineVerbatimEnvironment{Code}{Verbatim}{fontsize=\footnotesize}
\tikzset{
  box/.style={draw=black!60,rounded corners=2pt,align=center,
    inner sep=5pt,font=\small,minimum height=0.9cm},
  flow/.style={-{Stealth[length=2mm]},thick,draw=black!70},
  note/.style={font=\footnotesize,align=center}
}

\begin{document}
\begin{center}
  {\Large\bfseries HỆ THỐNG PHÂN TÍCH VIDEO NHIỀU CAMERA}\\[3pt]
  {\normalsize Workflow, xử lý sự kiện, vận hành và đánh giá hiệu năng}\\[2pt]
  {\small Báo cáo: 18/10/2026 \quad|\quad Benchmark local: 17/10/2026}
\end{center}
\vspace{-5pt}

\section{Workflow và trách nhiệm từng thành phần}
Hệ thống nhận nhiều luồng video, dùng YOLO và ByteTrack để nhận diện/theo dõi,
đếm vượt vạch và phát hiện vào vùng cấm. Kết quả được lưu SQLite, vẽ lên video
và phát RTSP; FastAPI cung cấp trạng thái, điều khiển camera và lịch sử sự kiện.

\begin{center}
\begin{tikzpicture}
\node[box,text width=2.65cm] (cfg) at (0,0)
  {\textbf{YAML}\\Camera, model, logging};
\node[box,text width=3.35cm] (mgr) at (4.1,0)
  {\textbf{SystemManager}\\Vòng điều phối trung tâm};
\node[box,text width=3.4cm] (api) at (8.6,0)
  {\textbf{FastAPI}\\Trạng thái, bật/tắt, lịch sử};
\draw[flow] (cfg)--(mgr);
\draw[flow,<->] (mgr)--(api);
\node[box,text width=3.05cm] (cap) at (-0.4,-1.75)
  {\textbf{CameraWorker $\times N$}\\FramePacket\\Frame mới nhất};
\node[box,text width=3.05cm] (infer) at (3.45,-1.75)
  {\textbf{InferenceWorker}\\DetectionPacket\\Bbox, class, confidence};
\node[box,text width=3.05cm] (track) at (7.3,-1.75)
  {\textbf{Tracker}\\TrackingPacket\\Bbox và track ID};
\node[box,text width=3.05cm] (event) at (11.15,-1.75)
  {\textbf{EventEngine}\\EventResults\\Crossing, zone entry};
\draw[flow] (mgr.south)--(infer.north);
\draw[flow] (cap)--(infer);
\draw[flow] (infer)--(track);
\draw[flow] (track)--(event);
\node[box,text width=4.05cm] (writer) at (7.0,-3.65)
  {\textbf{WriterWorker $\to$ EventWriter}\\SQLite + snapshot JPEG Base64};
\node[box,text width=4.45cm] (out) at (12.0,-3.65)
  {\textbf{Drawing $\to$ StreamWorker}\\FFmpeg $\to$ MediaMTX $\to$ RTSP};
\draw[flow] (event.south)--++(0,-0.25)-|(writer.north);
\draw[flow] (event.south)--++(0,-0.25)-|(out.north);
\end{tikzpicture}

{\small Hình 1. Pipeline dữ liệu; Manager gọi inference, tracker và event engine.
Preview local dùng cùng ảnh đã vẽ với nhánh RTSP.}
\end{center}

\begin{tabularx}{\linewidth}{@{}p{3.05cm}p{3.3cm}Y@{}}
\toprule
\textbf{Thành phần} & \textbf{Phạm vi / thực thi} & \textbf{Nhiệm vụ}\\
\midrule
CameraWorker & Mỗi camera, 1 thread & Đọc nguồn, reconnect, timestamp,
generation; queue 1 frame, thay frame cũ khi đầy.\\
InferenceWorker & Chung hoặc riêng model & YOLO suy luận trên batch do Manager gom;
trả kết quả về đúng camera.\\
Tracker & ByteTrack riêng/camera & Giữ track ID độc lập; một registry quản lý,
chạy trong vòng Manager.\\
EventEngine & Trạng thái riêng/camera & Nhớ phía vạch, vị trí trước, trạng thái
trong vùng và tổng đếm; chỉ tạo khi có quy tắc.\\
WriterWorker & Một thread ghi DB chung & Queue 128, ghi giao dịch cho frame có
sự kiện; queue đầy báo lỗi, không âm thầm bỏ sự kiện.\\
StreamWorker & Mỗi camera có RTSP & Một thread + một FFmpeg/camera; queue 1 frame,
giới hạn FPS xuất và tự phục hồi pipe.\\
API và logging & Dùng chung hệ thống & API đưa lệnh về Manager; listener riêng
ghi log, không dùng \texttt{print}.\\
\bottomrule
\end{tabularx}

\subsection*{Hai phương án được benchmark}
\shared{}: một YOLO xử lý batch nhiều camera.
\percam{}: mỗi camera giữ một YOLO, các model được gọi \textbf{tuần tự}.
Collector lấy tối đa một frame/camera, theo round-robin, tối đa
\texttt{batch\_size}; do đó model riêng thực tế nhận batch 1 trong mỗi lượt.
Tracker và EventEngine vẫn riêng theo camera ở cả hai phương án.
Hệ thống dùng threading và subprocess FFmpeg; chưa chia inference thành nhiều
process hay nhiều GPU.

\newpage
\section{Từ track ID đến sự kiện}
\subsection*{Điểm đại diện và kiểm tra hình học}
Với bbox $(x_1,y_1,x_2,y_2)$, code dùng \textbf{tâm bbox}
$C=((x_1+x_2)/2,(y_1+y_2)/2)$.
Quy tắc áp dụng cho cùng track ID trong cùng camera, trên tọa độ frame gốc.

\begin{center}
\begin{tikzpicture}
\node[box,text width=8.2cm] (root) at (0,0)
  {\textbf{TrackingPacket: bbox + track ID}\\
  Tính tâm $C$; đọc trạng thái trước của ID; loại trạng thái quá TTL};
\node[box,text width=6.8cm] (line) at (-4.1,-1.65)
  {\textbf{Có line $AB$: kiểm tra vượt vạch}\\
  Tâm đổi phía: $s(C_{\rm prev})s(C)<0$\\
  Đoạn chuyển động $C_{\rm prev}C$ giao đoạn $AB$};
\node[box,text width=6.8cm] (zone) at (4.1,-1.65)
  {\textbf{Có polygon: kiểm tra vào vùng}\\
  Tâm trong hoặc trên biên: \texttt{inside=True}\\
  Vào mới = inside hiện tại và trước đó ở ngoài};
\draw[flow] (root.south)--++(0,-0.25)-|(line.north);
\draw[flow] (root.south)--++(0,-0.25)-|(zone.north);
\node[box,text width=6.8cm] (count) at (-4.1,-3.2)
  {\textbf{Ghi một crossing và tăng bộ đếm chiều}\\
  $s(C)>0$: \texttt{left\_to\_right}; $s(C)<0$: chiều ngược\\
  Lưu phía/vị trí khác 0 gần nhất};
\node[box,text width=6.8cm] (entry) at (4.1,-3.2)
  {\textbf{Ghi một \texttt{zone\_entry}}\\
  Ở trong liên tục: không lặp sự kiện\\
  Ra rồi vào lại: được ghi lượt mới};
\draw[flow] (line)--(count);
\draw[flow] (zone)--(entry);
\end{tikzpicture}

{\small Hình 2. Hai nhánh độc lập; thiếu line hoặc polygon thì bỏ nhánh tương ứng.}
\end{center}

Tích có hướng của điểm $P$ đối với vạch có thứ tự $A\to B$:
\[
s(P)=(B_x-A_x)(P_y-A_y)-(B_y-A_y)(P_x-A_x).
\]
Sau khi xác nhận đổi phía, code kiểm tra thêm
$m(A)m(B)\leq0$, với $m$ là cùng phép tính trên đường $C_{\rm prev}\to C$.
Hai điều kiện giúp loại trường hợp chỉ cắt phần kéo dài ngoài hai đầu vạch.

\begin{center}
\begin{minipage}[t]{0.49\linewidth}
\centering
\begin{tikzpicture}[font=\footnotesize,scale=0.95]
\draw[flow] (0,1)--(4.7,1);
\node[above] at (0,1) {$A$}; \node[above] at (4.7,1) {$B$};
\draw[dashed,black!45] (4.7,1)--(6.15,1);
\draw[flow,draw=green!45!black] (2.0,2)--(2.0,0);
\fill[green!45!black] (2,2) circle(2pt);
\fill[green!45!black] (2,0) circle(2pt);
\node[left] at (2,2) {$C_{\rm prev}$};
\node[left] at (2,0) {$C$};
\draw[flow,dashed,draw=red!65!black] (5.6,2)--(5.6,0);
\node[below,align=center] at (2,-0.25) {Cắt đoạn: đếm 1 lượt};
\node[below,align=center] at (5.5,-0.25) {Ngoài đoạn:\\không đếm};
\end{tikzpicture}
\end{minipage}\hfill
\begin{minipage}[t]{0.49\linewidth}
\centering
\begin{tikzpicture}[font=\footnotesize,scale=0.95]
\filldraw[fill=blue!5,draw=black!60] (0.5,0)--(4.7,0)--(5,1.75)--(1,2)--cycle;
\draw[flow] (-0.1,0.8)--(2,0.8);
\draw[flow] (2,0.8)--(4,1.1);
\fill (-0.1,0.8) circle(2pt);
\fill[blue!70!black] (2,0.8) circle(2pt);
\fill[blue!70!black] (4,1.1) circle(2pt);
\node[above] at (-0.1,0.8) {Ngoài};
\node[above] at (2,0.8) {Trong};
\node[above] at (4,1.1) {Trong};
\node[below] at (2,-0.25) {Ngoài $\to$ trong: 1 entry};
\node[below] at (2,-0.65) {Trong $\to$ trong: không ghi lặp};
\end{tikzpicture}
\end{minipage}

{\small Hình 3. Minh họa crossing hợp lệ và chuyển trạng thái vào vùng.}
\end{center}

\begin{itemize}
\item \textbf{Chiều đếm:} tên \texttt{left\_to\_right}/\texttt{right\_to\_left}
phụ thuộc thứ tự hai đầu vạch, không nhất thiết là trái/phải trên màn hình.
Nếu tâm đúng trên vạch ($s=0$), giữ phía khác 0 trước đó để xét lượt sau.
\item \textbf{Vùng cấm:} lần đầu thấy một ID đã ở trong vùng cũng tạo entry;
biên polygon tính là trong. Một ID có thể được đếm nhiều lần khi vượt/ra vào lại.
\item \textbf{Quản lý trạng thái:} ID không thấy quá 60 giây bị xóa trạng thái
(TTL mặc định trong schema); reconnect, gap lớn hoặc bật/tắt sẽ reset danh tính,
giữ tổng đếm trong phiên. Đổi ID có thể gây đếm lại đối tượng.
\item \textbf{Giới hạn hiện tại:} chưa có vùng đệm, cooldown hay xác nhận ba
quan sát cho sự kiện. Thiết lập hai frame xác nhận của ByteTrack là logic
tracker, không phải bộ lọc sự kiện.
\end{itemize}

Writer chỉ ghi các \texttt{zone\_entry} hoặc
\texttt{crossed\_left\_to\_right}/\texttt{crossed\_right\_to\_left}.
Mỗi bản ghi gồm thời điểm \texttt{datetime.now()}, camera, track ID, loại sự
kiện và snapshot JPEG chất lượng 80 lưu Base64. Các sự kiện cùng frame ghi
trong một transaction SQLite (WAL); API đọc bằng connection riêng.

\newpage
\section{Cấu hình và vận hành}
\subsection*{Ba nhóm cấu hình}
\begin{tabularx}{\linewidth}{@{}p{3.2cm}Y@{}}
\toprule
\textbf{File / thành phần} & \textbf{Thiết lập hiện tại và ý nghĩa}\\
\midrule
\texttt{configs/model.yaml} &
Weights \texttt{best.pt} tại
\texttt{yolo/results/detect/runs/train/person\_vehicle/weights/};
\texttt{device=cuda:0}, \texttt{imgsz=640}, \texttt{confidence=0.5},
\texttt{classes=[0,\ldots,9]}. Ngưỡng confidence lọc detection;
tên class theo model đã huấn luyện.\\
Runtime trong model YAML &
\texttt{model\_mode=shared}, \texttt{batch\_size=8},
\texttt{tracking\_fps=15}, \texttt{reset\_after\_seconds=3};
queue sự kiện 128; DB \texttt{data/events.sqlite3}.
Batch 8 là tối đa, không bắt buộc đợi đủ 8 frame.\\
ByteTrack trong source &
Activation 0,5; high-confidence 0,6; lost buffer 30;
minimum consecutive frames 2; minimum IoU 0,1.
\texttt{tracking\_fps=15} là tham số tracker, không khóa FPS xử lý ở 15.\\
\texttt{configs/cameras.yaml} &
ID duy nhất, URL nguồn, URL output riêng, tọa độ, line/polygon tùy chọn;
không có cả hai vẫn chạy detect/tracking/video.\\
\texttt{configs/logging.yaml} &
INFO; console + UTF-8 \texttt{logs/system.log}; xoay file 10 MiB,
giữ 5 bản cũ; queue logging 4096; status mỗi 30 giây.\\
\bottomrule
\end{tabularx}

Ví dụ một camera (các điểm chỉ minh họa, phải khớp cảnh thực tế):
\begin{Code}
cameras:
  - id: "cam_01"
    coordinates: "normalized"
    source: {type: "rtsp", url: "rtsp://localhost:8554/live1.sdp"}
    output:
      rtsp_url: "rtsp://localhost:8554/output/cam_01"
      width: 800
      height: 450
      fps: 15
    lines:
      point: [[0.1, 0.5], [0.9, 0.5]]
    prohibit_area:
      polygon: [[0.2, 0.2], [0.8, 0.2], [0.8, 0.8], [0.2, 0.8]]
\end{Code}
\textbf{Tọa độ:} \texttt{normalized} là $[0,1]$, đổi thành $(xW,yH)$;
\texttt{pixels} dùng pixel frame nguồn. Line có đúng hai đầu khác nhau;
polygon có ít nhất ba điểm, diện tích khác 0. Preview/RTSP resize sau xử lý:
ảnh chụp display khác kích thước cần quy đổi về frame gốc trước khi đưa vào YAML.
Chẳng hạn ảnh 1920$\times$1080 của frame 1280$\times$720 cần nhân tọa độ với $2/3$
nếu chỉ resize toàn ảnh, không crop/padding.

\subsection*{Chạy local và sử dụng API}
Cần môi trường Python có dependencies trong \texttt{requirements.txt},
PyTorch phù hợp CUDA, weights và FFmpeg trong PATH; nguồn RTSP đã được phát.
Từ thư mục project, một lệnh chạy API cùng preview:
\begin{Code}
python -m src.manager --display --model_mode shared --log_file logs/local.log
\end{Code}
Bỏ \texttt{--display} để chạy không cửa sổ; \texttt{--no\_stream} tắt phát RTSP;
\texttt{--cameras}/\texttt{--model} đổi YAML; dùng
\texttt{--model\_mode per\_camera} cho đối chứng.
API mặc định tại \texttt{127.0.0.1:8000}, Swagger \texttt{/docs};
chạy đúng một Uvicorn worker để tránh tạo trùng pipeline/model.

\begin{tabularx}{\linewidth}{@{}p{6cm}Y@{}}
\toprule
\textbf{Endpoint} & \textbf{Mục đích}\\
\midrule
\texttt{GET /status} & Toàn pipeline, queue và metrics từng camera.\\
\texttt{GET /cameras}, \texttt{GET /cameras/\{id\}} & Trạng thái,
enabled, FPS và lỗi.\\
\texttt{PATCH /cameras/\{id\}} & Body \texttt{\{"enabled": false\}} hoặc
\texttt{true}; điều khiển camera đã có trong YAML.\\
\texttt{GET /events} & Lọc \texttt{camera\_id}, \texttt{start\_time},
\texttt{end\_time}; phân trang \texttt{limit} (100, tối đa 500), \texttt{offset}.\\
\bottomrule
\end{tabularx}
Thời gian truy vấn dạng \texttt{2026-10-17T16:00:00}, giờ máy không kèm múi giờ;
bao gồm cả hai mốc. Bật/tắt không ghi lại YAML, giữ model và tổng đếm.
Ctrl+C hoặc Q/Esc/đóng preview sẽ dừng pipeline/API và xả queue sự kiện.

\newpage
\section{Phục hồi, logging và triển khai Docker/MediaMTX}
\subsection*{Kết nối lại và theo dõi vận hành}
\textbf{Nguồn camera:} RTSP có timeout mở/đọc 3 giây. Lỗi mở hoặc đọc
$\to$ disconnected $\to$ xóa frame chờ, đóng nguồn $\to$ đợi 2 giây rồi thử lại.
Mỗi lần mở thành công tăng \texttt{generation}; chỉ báo connected sau frame tốt
đầu tiên. \texttt{connection\_attempts} đếm mọi lần mở,
\texttt{reconnect\_attempts} đếm lần thử sau khi từng có hình;
\texttt{reconnect\_count} chỉ tăng khi thật sự nhận lại frame.
Manager bỏ frame cũ hơn 3 giây; generation đổi hoặc gap $>3$ giây sẽ reset
tracker/trạng thái sự kiện/FPS, giữ tổng đếm hai chiều.

\textbf{Đầu ra:} FFmpeg lỗi/broken pipe được khởi động lại; lần gửi lỗi đợi
2 giây trước khi thử tiếp. Watchdog xử lý pipe ghi treo quá 3 giây.
Encoder là \texttt{libx264} trên CPU, preset \texttt{ultrafast},
tune \texttt{zerolatency}, RTSP qua TCP; chưa dùng NVENC.

\textbf{Logging:} dùng queue/listener để ghi console và file, có timestamp,
level, module, thread, camera ID; che thông tin đăng nhập trong URL.
INFO ghi khởi động, kết nối/phục hồi, commit DB và status; DEBUG ghi chi tiết
sự kiện/inference; WARNING/ERROR ghi mất nguồn, lỗi FFmpeg/worker/DB.
\texttt{--log\_level DEBUG} tăng chi tiết, \texttt{--log\_file} đổi nơi lưu.
Status mỗi 30 giây gồm FPS capture/xử lý/output, tổng frame, reconnect, queue,
drop, send failure và độ trễ frame-age/inference/processing/queue/send/E2E.
FPS và latency last/avg/P95 phản ánh \textbf{cửa sổ 5 giây}; không phải thống kê
toàn phiên. FPS xử lý cũng được vẽ lên từng camera.

\subsection*{Vai trò MediaMTX và mạng triển khai}
MediaMTX là \textbf{server luồng độc lập}; ứng dụng không tự tạo server RTSP.
Với video file, FFmpeg phát nguồn lên server trước; với camera thật, có thể
đọc URL camera trực tiếp. Nhánh output publish sang đường dẫn riêng.
\begin{center}
\begin{tikzpicture}
\node[box,text width=3.3cm] (s) at (0,0)
  {Camera / FFmpeg nguồn\\\texttt{/live1.sdp}};
\node[box,text width=3.3cm] (m) at (4.4,0)
  {\textbf{MediaMTX:8554}\\Nhận và phân phối RTSP};
\node[box,text width=4.8cm] (a) at (10,0)
  {\textbf{Ứng dụng Docker + GPU}\\Đọc nguồn; xử lý; FFmpeg xuất};
\draw[flow] (s)--(m);
\draw[flow] (m.north east) to[bend left=16]
  node[above,font=\footnotesize]{đọc nguồn} (a.north west);
\draw[flow] (a.south west) to[bend left=16]
  node[below,font=\footnotesize]{publish \texttt{/output/cam\_01}} (m.south east);
\end{tikzpicture}
\end{center}
Client/VLC đọc \texttt{rtsp://127.0.0.1:8554/output/cam\_01} trên host.
Docker Desktop dùng \texttt{host.docker.internal:8554} để gọi MediaMTX trên host,
như YAML Docker hiện tại; \texttt{localhost} trong container là chính container.
Nếu cả hai chạy cùng Docker network trên server, đổi host RTSP sang tên service
MediaMTX. MediaMTX có thể chạy binary hoặc container; với bridge network, dùng
RTSP/TCP và publish cổng 8554.\footnote{Tài liệu:
\href{https://mediamtx.org/docs/kickoff/install}{MediaMTX -- Install};
\href{https://docs.docker.com/desktop/features/networking/networking-how-tos/}
{Docker Desktop -- Networking how-tos}.}

\subsection*{Lệnh mẫu trên Windows CMD}
Dùng một MediaMTX đang chạy sẵn; nếu chưa có server, có thể khởi tạo bằng
lệnh đầu. Các bind mount nguồn phải tồn tại. Dockerfile đóng gói môi trường
PyTorch/CUDA, dependencies, FFmpeg và source; weights, dữ liệu, log được mount.
Base image đang khai báo: PyTorch 2.14.1, CUDA 12.6, cuDNN 9 runtime.
\begin{Code}
docker run -d --name mediamtx -e MTX_RTSPTRANSPORTS=tcp -p 8554:8554 bluenviron/mediamtx:1
docker build -t multi-camera:dev .
set "WEIGHTS_DIR=%CD%\yolo\results\detect\runs\train\person_vehicle\weights"
set "WEIGHTS_TARGET=/app/yolo/results/detect/runs/train/person_vehicle/weights"
mkdir logs
docker run -d --name multi-camera --gpus all --stop-timeout 90 ^
  -p 127.0.0.1:8000:8000 ^
  --mount "type=bind,source=%CD%\configs,target=/app/configs,readonly" ^
  --mount "type=bind,source=%WEIGHTS_DIR%,target=%WEIGHTS_TARGET%,readonly" ^
  --mount "type=volume,source=multi-camera-data,target=/app/data" ^
  --mount "type=bind,source=%CD%\logs,target=/app/logs" ^
  multi-camera:dev python -m src.manager --host 0.0.0.0 --port 8000 ^
  --cameras /app/configs/cameras.docker.yaml --model_mode shared ^
  --log_file /app/logs/docker.log
\end{Code}
Container chạy API và RTSP không mở preview OpenCV; host xem bằng client RTSP.
Log xuất trực tiếp ra \texttt{logs/docker.log}, DB nằm trong volume.
\textbf{Lưu ý source hiện tại:} Dockerfile còn CMD gọi \texttt{src.api}, trong khi
entry point đã chuyển sang \texttt{src.manager}; lệnh trên ghi đè CMD đúng.
Đây là hướng dẫn triển khai, không phải một lượt benchmark Docker đã đo.

\newpage
\section{Môi trường đo, benchmark và lựa chọn}
\subsection*{Môi trường và phương pháp}
\textbf{Máy đo do người thực hiện xác nhận:} AMD Ryzen 7 7735HS,
NVIDIA RTX 4050 Laptop 6 GB VRAM, RAM 16 GB.
Chạy local trên Windows; sáu log tại
\texttt{benchmark/local/\{1cam,4cam,8cam\}/\{share,per\}.log}.
Log xác nhận CUDA0, batch tối đa 8, display và streaming bật,
output $800\times450$, giới hạn 15 FPS.
Thông số weights/confidence/classes/imgsz trình bày ở mục 3 là
\textbf{cấu hình source hiện tại}; log chưa lưu đầy đủ để truy nguyên chính
xác mọi tham số/phiên bản của các phiên đo. Docker chưa có số liệu đối chứng.

Mỗi cấu hình có một phiên thành công. Chọn snapshot chạy định kỳ thứ 3 đến 9,
khoảng giây 90--272, được 180,04--180,79 giây/cấu hình, tất cả camera connected.
Loại phần khởi động và tiền tố lần chạy lỗi trong \texttt{8cam/per.log}.
Thông lượng dùng toàn khoảng; latency lấy sáu snapshot sau mốc đầu:
\[
 F_{\rm total}=\frac{\Delta\texttt{processed\_frames}}{\Delta t},
 \qquad F_{\rm cam}=\frac{F_{\rm total}}{N_{\rm camera}}.
\]
\vspace{-8pt}
Ví dụ 4 camera \shared{}: $9\,374/180{,}146=52{,}04$ FPS.
Latency là trung bình các \texttt{avg\_ms} cửa sổ 5 giây, gộp camera đồng trọng
số; không phải trung bình toàn bộ frame trong 180 giây.

\subsection*{Số liệu tổng hợp}
\begin{center}
\small
\setlength{\tabcolsep}{7pt}
\begin{tabular}{@{}crrrrrr@{}}
\toprule
\shortstack{Số\\camera} & \multicolumn{1}{c}{Phương án} &
\shortstack{FPS\\tổng} & \shortstack{FPS/cam\\TB} &
\shortstack{$L_{\rm proc}$\\(ms)} &
\shortstack{$L_{\rm pipe}$\\(ms)} & \shortstack{FPS xuất\\TB/cam}\\
\midrule
1 & \shared & 28,61 & 28,61 & 37,46 & 61,28 & 12,71\\
1 & \percam & 28,58 & 28,58 & 37,71 & 63,82 & 12,72\\
\addlinespace
4 & \shared & \textbf{52,04} & \textbf{13,01} & 71,33 & 113,74 & 12,77\\
4 & \percam & 36,16 & 9,04 & 102,58 & 107,30 & 9,04\\
\addlinespace
8 & \shared & \textbf{29,26} & \textbf{3,66} & 212,64 & 220,84 & 3,66\\
8 & \percam & 24,61 & 3,08 & 286,88 & 293,27 & 3,08\\
\bottomrule
\end{tabular}
\end{center}

$L_{\rm proc}$: từ khi Python nhận frame đến sau tracking/event và enqueue
writer, trước drawing. $L_{\rm pipe}$: trường \texttt{e2e\_ms}, từ lúc nhận frame
đến write/flush pipe FFmpeg; chưa gồm hoàn tất encode, mạng, RTSP server/client.
FPS xuất dùng $\Delta\texttt{sent\_frames}/\Delta t$.

\subsection*{Nhận xét trên số liệu}
\begin{itemize}
\item \textbf{1 camera:} chênh lệch thông lượng chỉ 0,09\%, hai mode gần tương
đương; một lượt đo chưa chứng minh được ưu thế đáng kể.
\item \textbf{4 camera:} \shared{} tăng thông lượng \textbf{43,9\%}, giảm
$L_{\rm proc}$ \textbf{30,5\%}; từng camera đạt 12,92--13,05 FPS so với
9,02--9,05. Tuy nhiên $L_{\rm pipe}$ \textbf{tăng 6,0\%}; log ghi nhận chờ queue
đầu ra 36,89 so với 0,33 ms, và 169 frame output bị thay thế so với 0.
Vì vậy ưu thế xử lý không đồng nghĩa mọi nhánh đều ít trễ hơn.
\item \textbf{8 camera:} \shared{} tăng thông lượng \textbf{18,9\%}, giảm
$L_{\rm proc}$ \textbf{25,9\%} và $L_{\rm pipe}$ \textbf{24,7\%}.
Từng camera đạt 3,40--3,86 FPS. Nếu cần 15 FPS xử lý/camera, cả cấu hình 4 và 8
camera hiện chưa đạt; chưa đủ dữ liệu tách riêng nút thắt GPU/decode/tracking/encode.
\item \textbf{Ổn định trong khoảng đo:} không tăng reconnect/send-failure;
queue sự kiện bằng 0 tại các mốc lấy mẫu. Với một camera, hai mode có khoảng
2\,860 frame output bị thay thế do tốc độ xử lý lớn hơn tốc độ xuất.
Đây là drop ở nhánh streaming, không phải số sự kiện bị bỏ sót.
\end{itemize}

\textbf{Lựa chọn: dùng \shared{} làm mặc định khi chạy nhiều camera.}
Phương án này đạt thông lượng cao hơn, trễ xử lý thấp hơn và chỉ cần một model.
Giữ \percam{} để đối chứng; chưa có lý do hiệu năng từ các log này để chọn mặc
định. Tăng số camera vẫn cần đo lại theo FPS và độ trễ yêu cầu.

\subsection*{Giới hạn và hướng kiểm chứng}
Chỉ một lần chạy/cấu hình; chưa có khoảng tin cậy, số đo CPU/GPU/RAM/VRAM
theo thời gian hay nhãn đối chứng độ chính xác sự kiện. Nguồn có capture FPS
khác nhau (khoảng 24/30/60/120), nên ưu tiên so sánh hai mode cùng số camera;
không suy ra khả năng scale tuyến tính. Hai mode đồng thời khác số bản model
và batching (batch nhiều camera so với tuần tự batch 1); kết luận áp dụng cho
hai cách triển khai hiện tại, chưa tách riêng tác động của từng yếu tố.
Cần lặp ít nhất ba lượt/mode trên cùng
video, model và cấu hình; đo thêm tài nguyên, tỷ lệ đếm đúng và điều kiện mất
kết nối. So local/Docker phải giữ display, output, encoder và nguồn tương
đương; dữ liệu hiện tại chưa cho phép kết luận Docker nhanh hơn.

% Nguồn kiểm chứng: src/manager.py, camera_worker.py, event_engine.py,
% tracker.py, output_workers.py, stream_ouput.py, event_writer.py, api.py;
% configs/model.yaml, cameras.yaml, cameras.docker.yaml, logging.yaml.
% FPS: snapshot định kỳ thứ 3 -> thứ 9, 17/10/2026 (+07:00).
% 1cam/per.log:   16:21:35.354 -> 16:24:35.406; 5146 / 180.052 s.
% 1cam/share.log: 16:27:48.930 -> 16:30:48.966; 5150 / 180.036 s.
% 4cam/per.log:   16:40:44.899 -> 16:43:45.192; 6519 / 180.293 s.
% 4cam/share.log: 16:46:11.942 -> 16:49:12.088; 9374 / 180.146 s.
% 8cam/per.log:   17:05:31.300 -> 17:08:32.089; 4450 / 180.789 s.
% 8cam/share.log: 17:14:23.936 -> 17:17:24.402; 5281 / 180.466 s.
% Latency: seq[3:9] của từng camera, loại snapshot khi dừng.
\end{document}

